\chapter{Methodology}
\section{Research Design and Conceptual Framework}\label{sec:research-design}
This applied quantitative study evaluates a text-only, three-class complaint classifier using a fixed synthetic corpus. Phase I concerns corpus documentation, sparse feature extraction, model selection and offline classification evaluation. Phase II specifies a proposed human-governed reporting system. The two phases are related but have different evidence requirements: offline model scores do not establish that resident accounts, administrator review, accepted-case management or feedback-driven retraining work in practice. Figure~\ref{fig:overall} separates these boundaries.
\begin{figure}[H]\centering\includegraphics[width=.96\textwidth,height=.78\textheight,keepaspectratio]{figures/conceptual_framework.pdf}\caption{Two-phase conceptual framework: measured classification research and proposed operational workflows. Human feedback enters authorized offline evaluation rather than changing a live model automatically.}\label{fig:overall}\end{figure}

\section{Dataset Provenance, Labeling and Integrity}\label{sec:dataset-protocol}
The canonical research input is \path{data/Barangay_Complaint_Dataset_24000.csv}. Its 24,000 narratives are researcher-constructed synthetic examples, not authentic resident archives. The source register, category register and labeling guidelines document construction choices and source-to-category associations. Legislation and official reporting practices establish contextual relevance; they do not certify individual generated incidents or natural class prevalence \cite{ra7160,dcr,ra9003,deleon,quimbo,porlas,long,chim}. The corpus uses 60 Valid topics, 28 Suspicious review reasons and 30 Spam types as 118 coverage categories, not 118 prediction classes. Generation seed 20261026 and the original source bytes are preserved; reproduction never requires regenerating the corpus.

\begin{table}[H]\centering\caption{Fixed corpus composition and original partitions}\label{tab:corpus-final}\small\begin{tabularx}{\textwidth}{l r X}\toprule Dimension & Records & Interpretation\\\midrule
Valid / Suspicious / Spam & 8,000 each & Deliberately balanced supervisory labels\\
Bisaya / Filipino / English / Mixed & 6,000 each & Assigned, independently unverified language groups\\
Training & 19,213 & Feature/model fitting and grouped resampling\\
Validation & 2,399 & Existing development evaluation\\
Synthetic test & 2,388 & Original test partition, already evaluated\\\bottomrule\end{tabularx}\end{table}

\texttt{complaint\_text} is the only predictive input; \texttt{classification} is the target. Category descriptions, language tags, label rules, split flags, source references and family identifiers are audit metadata and are excluded from model features. The archived \path{data/ALL_SYNTHETIC_V2.csv} is a byte-identical export of the canonical CSV: both have SHA-256 \path{35414af2dc63e2870928bd1853c557ed301dbdae80416a0bf37be00648d92e28}. All 21 fields and record order match, including narratives, labels, language assignments, categories, family IDs and partitions. Split exports are checked against their corresponding canonical records before use. The standalone Overleaf project includes this verified mirror for self-contained archival documentation; the complete research repository identifies one canonical input path.

The text-label rubric is independent of factual truth, jurisdiction and administrative disposition. \textbf{Valid} denotes an interpretable community concern or request for action; a short, emotional or anonymous report can qualify. \textbf{Suspicious} denotes report-like text whose essential meaning is missing, contradictory or unresolved; it is not an allegation of dishonesty. \textbf{Spam} denotes irrelevant, promotional, deceptive or nonsensical non-complaint text. A clear report \emph{about} a scam can be Valid. Administrative acceptance does not retroactively convert a Suspicious original into a Valid training example; a clarified narrative is a separate version to label independently \cite{rottger2022,fleisig2024}.

The audit checks unique record IDs, nonempty texts, label/partition counts, duplicate texts and cross-partition scenario-family separation. Unicode-normalized exact matching and nearest-training TF--IDF cosine similarity provide supplementary diagnostics; similarity is not an adjudicated duplicate definition or a basis for deleting records. Generation and partitioning are illustrated in Figure~\ref{fig:dataset}.
\begin{figure}[H]\centering\includegraphics[width=.94\textwidth,height=.72\textheight,keepaspectratio]{figures/dataset_pipeline.pdf}\caption{Documented scenario construction and audit process for the preserved synthetic corpus. Category coverage and source traceability do not establish authentic population frequencies.}\label{fig:dataset}\end{figure}

All records retain \texttt{language\_verified=No} and \texttt{human\_label\_reviewed=No}. No inter-rater agreement study was completed. Independent blinded annotation, adjudication and fluent language review remain necessary before stronger validity claims. Shared generator rules, lexical patterns, class-dependent lengths and unknown real class prevalence remain threats even with nonoverlapping family IDs \cite{artstein,benderfriedman,datasheets,gururangan,shaib}. Chapter~4 reports integrity and shortcut findings; these controls do not establish authentic-report generalization.

Descriptive narrative-length analysis is retained as a dataset-quality diagnostic in Figure~\ref{fig:m1lengths}. The distributions help identify superficial class cues that grouped classification evaluation alone cannot exclude; they are properties of the generated corpus rather than estimates of resident writing behavior.
\begin{figure}[H]\centering\includegraphics[width=.92\textwidth,height=.40\textheight,keepaspectratio]{figures/m1_length_violin.pdf}\caption{Preserved descriptive narrative-length distributions in the synthetic corpus. Class-dependent surface properties motivate the shortcut diagnostics reported in Chapter~4.}\label{fig:m1lengths}\end{figure}

\section{Hardware and Software Environment}\label{sec:environment-final}
Table~\ref{tab:environment-final} distinguishes specifications in the supplied Milestone~2 reference from the subsequent verification environment. The reference document explicitly records the CPU/GPU, two 8 GB Samsung DDR5-5600 modules and historical OS/software versions. A read-only device inventory on 27 October 2026 corroborated the CPU, GPUs, memory modules and Windows edition/version; the current Windows patch build differs. A later inventory cannot prove the exact operating-system patch or installed libraries at the original training time. The supplied historical software record and later captured package manifest are therefore reported separately.
\begin{table}[H]\centering\caption{Recorded research environment and later verification environment}\label{tab:environment-final}\small\begin{tabularx}{\textwidth}{p{.20\textwidth}XX}\toprule Component & Supplied original record & Later verification record\\\midrule
CPU & Intel Core i7-13700HX & Same CPU model, device query\\
GPU & NVIDIA RTX 4060 Laptop; Intel UHD & Same GPU models, device query\\
RAM & 16 GB; two 8 GB Samsung DDR5-5600 & Two 8 GB Samsung modules at 5600, device query\\
OS & Windows 11 Home 25H2, build 26200.9457 & Windows 11 Home 25H2, build 26200.9711\\
Python & 3.13.5 & 3.13.5\\
scikit-learn & 1.8.0 & 1.8.0\\
NumPy / SciPy & 2.3.5 / 1.17.0 & 2.4.4 / 1.17.1\\
pandas / joblib & 2.2.3 / 1.5.3 & 3.0.2 / 1.5.3\\
Plot/QA dependencies & Not fully captured in original record & Matplotlib 3.11.2; Pillow 12.3.0; threadpoolctl 3.6.0\\\bottomrule\end{tabularx}\end{table}

The actual scikit-learn TF--IDF, LinearSVC and Naive Bayes implementations are CPU-based. No CUDA/GPU estimator or GPU acceleration was used by these scripts; presence of an NVIDIA GPU is an inventory fact. Later verification establishes numerical repeatability of recorded outputs, not bitwise reconstruction of the original environment. The pinned research dependency file represents the verified environment; the supplied historical versions remain identifiable in archival records. Next.js/TypeScript, FastAPI, PostgreSQL and their planned migration/authentication dependencies belong to the proposed application and were not dependencies of the reported offline experiments. pdfLaTeX and the included packages compile the manuscript; their compile logs are artifact evidence, not ML training records.

\section{Reference Pipelines and Baseline Selection}\label{sec:baseline-final}
Multinomial Naive Bayes is the standard sparse-text baseline. It estimates class priors and class-conditional feature likelihoods with additive smoothing, using a conditional-independence approximation. With sparse nonnegative TF--IDF features, fitting primarily accumulates class-wise feature statistics and inference evaluates class scores; this makes it a computationally inexpensive and useful reference for bag-of-words classification. The assumption is imperfect for correlated words, but its simplicity tests whether the proposed SVM adds useful discrimination on the defined complaint labels. LinearSVC instead learns regularized discriminative decision boundaries in the sparse feature space \cite{sklearn,sklearntfidf}.

\begin{table}[H]\centering\caption{Recorded reference pipeline configurations}\label{tab:reference-final}\small\begin{tabularx}{\textwidth}{p{.27\textwidth}XX}\toprule Component & Linear SVM reference & Naive Bayes baseline\\\midrule
Input & Raw complaint narrative & Raw complaint narrative\\
Word n-grams & Selected $(1,1)$ & Fixed $(1,2)$\\
Minimum document frequency & Selected 2 & Fixed 2\\
Maximum document frequency & Fixed 0.95 & Fixed 0.95\\
TF--IDF weighting & Sublinear TF; smooth IDF; L2 & Sublinear TF; smooth IDF; L2\\
Estimator & LinearSVC; selected $C=0.1$ & MultinomialNB; $\alpha=1$\\
Randomness / iterations & Seed 20261026; max\_iter 5000 & Deterministic given input/features\\\bottomrule\end{tabularx}\end{table}

Both vectorizers use scikit-learn's default word analyzer, lowercasing and token pattern, with no stopword list. The reference scripts pass raw narratives directly to the saved vectorizer; the demo's separate URL/punctuation-stripping helper was not inserted into the historical study. The vectorizer and estimator form one Pipeline so vocabulary and IDF are fitted inside each training fold. Table~\ref{tab:reference-final} contrasts \emph{complete pipelines}: different n-gram ranges prevent attributing the entire difference solely to classifier choice. No historical settings were changed to produce matched features.

\section{Partitioning, Hyperparameter Search and Stability Analysis}\label{sec:tuning-final}
The stored \texttt{split} column fixes the original 19,213/2,399/2,388 membership. Scenario-family IDs remain disjoint across partitions and are passed as groups to cross-validation; neither validation nor test examples fit vocabulary, IDF, estimator parameters or probability calibration. This prevents direct family-ID overlap, although related wording may remain across different IDs.

The SVM uses exhaustive \texttt{GridSearchCV} over the 18 Cartesian-product configurations in Table~\ref{tab:search-final}. Three-fold \texttt{StratifiedGroupKFold} uses \texttt{n\_splits=3}, \texttt{shuffle=True} and \texttt{random\_state=20261026}, scoring candidates on training-only groups with Macro-F1. The selected estimator is refitted on the full training partition. Fixed vectorizer options and estimator defaults are preserved and exported in the saved configuration evidence. In a tie, the deterministic parameter enumeration determines the first selected configuration; the tie is not proof of intrinsic parameter superiority.
\begin{table}[H]\centering\caption{Original 18-configuration SVM tuning strategy}\label{tab:search-final}\small\begin{tabularx}{\textwidth}{p{.40\textwidth}X}\toprule Setting & Values or method\\\midrule
Variable LinearSVC $C$ & $0.1,1,10$\\
Variable word n-grams & $(1,1),(1,2)$\\
Variable \texttt{min\_df} & $2,5,10$\\
Fixed \texttt{max\_df} & 0.95\\
Fixed \texttt{sublinear\_tf} & True\\
Search and selection & GridSearchCV; three-fold scenario-grouped mean Macro-F1\\
Configurations / fitting policy & $3\times2\times3=18$; TF--IDF inside Pipeline\\
Seed / execution & 20261026; CPU; search n\_jobs=1\\\bottomrule\end{tabularx}\end{table}

The separate five-fold \texttt{StratifiedGroupKFold} procedure fits the selected SVM and fixed Naive Bayes afresh on training folds to describe stability. It is not another hyperparameter search, nested model selection or five independent datasets. Each fold checks group disjointness and refits TF--IDF. Fold means and sample standard deviations describe partition sensitivity, not independent population confidence intervals. The fixed validation set supports development diagnostics, while the original synthetic test has already been inspected and is not a new untouched test for future tuning. Figure~\ref{fig:evaluation} summarizes allocation.
\begin{figure}[H]\centering\includegraphics[width=.96\textwidth,height=.72\textheight,keepaspectratio]{figures/m2_unified_evaluation.pdf}\caption{Original experimental allocation: training-only three-fold search and separate five-fold stability analysis, fixed validation, and the already evaluated synthetic test.}\label{fig:evaluation}\end{figure}

\section{Evaluation Metrics and Error Analysis}\label{sec:metrics-final}
Let $K=3$ classes, $N$ evaluated narratives, $y_i$ the reference text label of record $i$, and $\hat y_i$ its prediction. For class $c$, $TP_c$ counts correctly predicted class-$c$ records, $FP_c$ counts other-class records predicted as $c$, $FN_c$ counts class-$c$ records predicted otherwise, and $TN_c$ counts records neither reference nor predicted as $c$. Let $n_c$ be reference support in class $c$ and $\mathbf{1}[A]$ equal one when condition $A$ is true, otherwise zero. For single-label multiclass classification:
\begin{align}
\mathrm{Accuracy}&=\frac{1}{N}\sum_{i=1}^{N}\mathbf{1}[y_i=\hat y_i]=\frac{\sum_{c=1}^{K}TP_c}{N},\\
P_c&=\frac{TP_c}{TP_c+FP_c}, & R_c&=\frac{TP_c}{TP_c+FN_c},\\
F_{1,c}&=\frac{2P_cR_c}{P_c+R_c}=\frac{2TP_c}{2TP_c+FP_c+FN_c},\\
P_{\mathrm{macro}}&=\frac{1}{K}\sum_c P_c,&R_{\mathrm{macro}}&=\frac{1}{K}\sum_c R_c,\\
F_{1,\mathrm{macro}}&=\frac{1}{K}\sum_c F_{1,c},&F_{1,\mathrm{weighted}}&=\frac{1}{N}\sum_c n_cF_{1,c}.
\end{align}
Zero-denominator class scores are set to zero explicitly. Accuracy expresses overall agreement; precision describes how often predictions of a class are correct; recall describes how many reference examples of that class are recovered. F1 balances precision and recall. Macro-F1 is the primary measure because it gives each of the three operational text labels equal influence rather than letting a common class dominate. Weighted-F1 reflects reference support. Macro-F1 is the mean of class F1 scores, not the F1 of macro-precision and macro-recall.

The confusion entry $M_{ab}=\sum_i\mathbf{1}[y_i=a\land\hat y_i=b]$ counts reference class $a$ predicted as class $b$; rows are reference labels and columns are predictions. Per-class precision/recall/F1/support, mistakes, language-group support and confusion counts accompany aggregate metrics in Chapter~4. A Valid narrative predicted as Spam is particularly important because inappropriate rejection or deprioritization could prevent access to assistance. The proposed routing therefore sends Spam to human review rather than rejecting it automatically. Label agreement neither verifies an incident's factual truth nor measures complaint-resolution effectiveness.

Training-size learning curves measure Macro-F1 as sample count changes. \texttt{LinearSVC} does not emit the reported epoch-loss history; the fifteen-pass SGD hinge-loss trace is a distinct auxiliary estimator and must remain labeled accordingly. Length-only and character-only ablations are diagnostic checks for generator shortcuts rather than deployed alternatives. The reported word/character candidate is separately described below.

\section{Paired Statistical Comparison}\label{sec:mcnemar-final}
McNemar's test compares binary correctness for two classifiers on the same records \cite{mcnemar_nist}. ID-aligned predictions form $a$ (both correct), $b$ (baseline correct/SVM wrong), $c$ (baseline wrong/SVM correct) and $d$ (both wrong), with $N=a+b+c+d$. The null hypothesis is equal paired error probabilities; the significance level is $\alpha=0.05$. The implementation uses the nonnegative continuity-corrected statistic and exact two-sided binomial form:
\begin{align}
\chi^2_{\mathrm{McN}}&=\frac{\max(0,|b-c|-1)^2}{b+c},\qquad b+c>0,\\
X&\sim\mathrm{Binomial}(b+c,\tfrac12),\\
p_{\mathrm{exact}}&=\min\!\left(1,2\Pr[X\leq\min(b,c)]\right).
\end{align}
Here $X$ counts one discordant direction under equal probabilities; with no discordance, the statistic is zero and $p=1$. A small p-value indicates that the observed directional imbalance is unlikely under the paired null assumptions; it is not the probability that a model is reliable or that the null hypothesis is true. The exact test is preferred for the small discordance counts. Generated records share scenario/construction dependencies, so nominal p-values may overstate inferential evidence; results are limited to these complete pipelines on this synthetic corpus. No paired t-test or Friedman test is added to an unsupported design. Chapter~4 provides actual paired tables, statistics, p-values and interpretation.

\section{Experimental Documentation and Reproducibility}\label{sec:reproducibility-final}
The public repository is \url{https://github.com/rencze/nlp-barangay-complaint-classification-hitl}. Anonymous web access and a clean clone succeeded on 27 October 2026; the observed public main revision was \texttt{13dadbc3f82a11bce0383171624a78d4c3ccefc8}. That checkout contains the earlier demo source, demo CSV and broad dependency constraint, but its README does not contain installation/minimal-example instructions and it does not yet contain the full corrected study package. This distinguishes verified accessibility from completeness of remote artifacts. The revised local README, exact scientific dependency file, executable example, study scripts, configurations, data hashes, predictions and figures are delivered with this manuscript. Authenticated GitHub write access was unavailable during the audit; remote publication is pending and is not claimed as completed.

From the corrected project root, create a Python 3.13 environment, install \path{requirements.txt}, then run \path{python src/minimal_example.py}. It prints narrative-to-label predictions from a small demonstration corpus without claiming a research evaluation. The minimal source and relative demo path exist and were exercised from a clean packaged checkout with the pinned dependencies. The public checkout's older demo also runs when invoked explicitly, although its README did not document that command. Detailed commands, install logs and corrected/public-checkout distinctions are in the supporting audit reports.

Research input verification uses \path{python src/verify_dataset.py}. The original procedures remain identifiable as \path{manuscript/source/reproducibility_scripts/evaluate_models.py} and \path{manuscript/source/validation_suite/run_grouped_validation.py}; they retain their original feature, seed, partition and search settings. The full README documents explicit reproduction commands, but rerunning full training is unnecessary for the present final audit. Read-only \path{python src/final_evidence_audit.py} compares all dataset fields, aligns saved predictions with labels, reloads trusted saved models, recomputes metrics/McNemar and verifies training-fitted IDF without fitting any classifier. No personal local path or credential is required.

\subsection{Reproducibility Verification and Execution History}\label{sec:verification-history}
The original archived procedures were executed again during the earlier documentation work; logs, environment manifests, historical output copies and recorded-precision comparisons preserve that history. Those executions supported reproducibility, not a new independent study. The final Milestone~2 audit checks existing evidence and makes documentation corrections without repeating full model training. Historical outputs, subsequent verification outputs and candidate development outputs are separate named artifacts. No dataset was regenerated, relabeled or repartitioned.

Reference model bundles \path{results/verified_20261026/svm_research_only.joblib} and \path{results/verified_20261026/nb_research_only.joblib} contain the matching vectorizer/classifier. \path{src/evaluate_research.py} reloads a trusted bundle for saved-reference evaluation; reference inference passes raw text to that bundle rather than re-fitting vocabulary or applying the demo helper. These research artifact-loading checks are implemented and verified; their integration into authenticated application submission remains planned. Arbitrary uploaded joblib/pickle files must not be loaded. Prediction exports identify records, partitions, labels and model names; immutable artifact SHA-256 manifests and configurations provide lineage. Final figures are regenerated from actual metrics/predictions as 360-DPI PNG and vector PDF/SVG and presented with academic captions/tables, while CSV/JSON remain audit formats.

\subsection{Separately Evaluated Combined-Feature Candidate}\label{sec:candidate-final}
The candidate is distinct from the original SVM: Unicode NFKC, casefolding and whitespace normalization preserve negations and accented characters; word $(1,2)$ TF--IDF (min\_df 2, max\_df .95) is concatenated with character-within-word $(3,5)$ TF--IDF (min\_df 3, maximum 45,000 features). Both branches use sublinear TF and default L2 normalization, equal weights and LinearSVC $C=.1$. One fixed candidate was evaluated; no new search was performed. A calibrated variant uses sigmoid \texttt{CalibratedClassifierCV} with an ensemble of three precomputed scenario-disjoint training folds, each refitting its complete feature pipeline. Calibration folds do not fit their corresponding base estimators. Existing development validation evaluates the candidates; the historical test was not used for them \cite{sklearncalib,sklearnpitfalls}.

Both variants tied the reference's validation Macro-F1, so no classification improvement is demonstrated. Log loss, multiclass Brier score and ten-bin expected calibration error describe synthetic development probability behavior. These are completed development diagnostics, not independently verified operational calibration. The reference SVM has no calibrated probabilities; routing remains class-based, not an untested probability threshold. Chapter~4 reports candidate findings separately; live calibration approval and external validation remain future work.

\section{Proposed System and Operational Boundaries}\label{sec:system-final}
The selected architecture uses Next.js/TypeScript, FastAPI/Python, PostgreSQL and a versioned scikit-learn text pipeline. Predictions, human-verified text labels, dispositions and accepted-case statuses are distinct. Predicted Valid proceeds to policy-based accepted-case management; predicted Suspicious and Spam proceed to administrator review. Administrator confirmation/correction must preserve narrative versions and must not automatically convert acceptance/rejection into labels. Feedback eligibility, authorized offline training, independent candidate evaluation and model-version approval precede any live replacement. These are implementation requirements, not effects measured by the classifier experiment. The existing SQLite/static-HTML prototype does not meet the target authentication, routing, state, confidence or feedback gates. The detailed roadmap identifies existing code and missing modules without asserting implementation completion.
\input{chapters/operational_workflows.tex}

\section{Ethics, Limitations and Next Validation}\label{sec:limits-final}
The synthetic corpus avoids claiming private resident data collection, but unreviewed language/labels and a common construction process limit validity. Corpus balance is an experimental control rather than population prevalence. Future external validation needs independently authored/annotated texts, blinded agreement/adjudication, behavior-sensitive tests, adequate subgroup support and privacy review. Operational evaluation separately requires accounts, object-scoped authorization, urgency handling, audit integrity, referral/deadline policies, resident notices, case transitions, feedback quality and model rollback tests \cite{artstein,ra10173,ribeiro,nist}. No live HITL improvement, production deployment or complaint-resolution benefit is inferred from synthetic classification scores. The Milestone~2 traceability appendix links each requirement to manuscript sections and verifiable repository artifacts and explicitly retains remote-publication limitations.
